% VERSION B -- applies under the rule in README.md whenever Version A does not
% (including the floor case A1 < 0.20). Written before ToolRepair-12 results; fill \TBD{} only.
\providecommand{\TBD}[1]{\textbf{[#1]}}

\subsection{RQ4: Useful behavior}
\textbf{Pilot.} The ledger pilot (Table~\ref{tab:pilot}) diagnosed three adapter defects and showed, on two tasks, that sandboxed behavioral evidence admits correct model repairs that a finite catalog rejects. It could not estimate usefulness, so we ran a prespecified study.

\textbf{ToolRepair-12 design.} Twelve tool families each plant one defect in the harness's executable tool. Each family supplies a specification, a trusted reference, and three disjoint input sets: visible task inputs (graded by exact reference output), operator-held validation inputs (authority only), and oracle inputs (independent oracle only). Five pinned models run three seeds per family in three arms: original \harness{} (A1), the contract with the catalog policy (A2), and the contract with sandboxed evidence (A3), with a frozen reference and one injected red-team proposal per stream. The protocol, analysis, decision rule, and implementation hashes were fixed before qualification; the family is the unit of inference.

\textbf{Usefulness is not established.} Mean per-family success is \TBD{A1} for original \harness{}, \TBD{A2} under the catalog, and \TBD{A3} under sandboxed evidence (frozen: \TBD{A0}). \TBD{Choose one, per the README rule:} \emph{(floor)} Even without the contract, current open models rarely repair a defective tool and use it correctly, so the non-inferiority test is dominated by a floor effect and cannot show whether the contract costs usefulness. \emph{(inferior)} The contract with sandboxed evidence falls short of original \harness{} by \TBD{A1-A3} points (95\% CI \TBD{lo, hi}), outside the prespecified 10-point margin. The catalog policy \TBD{does/does not} lose further success relative to sandboxed evidence (H1: \TBD{diff, CI}). Correct repairs were proposed in \TBD{x} streams and executed in \TBD{y}; the dominant failures were \TBD{e.g., re-emitting the defective tool, answering without the tool, protocol failures}.

\textbf{Safety holds independently of usefulness.} Original \harness{} executed \TBD{n} contract-violating proposals (\TBD{natural} natural, \TBD{injected}/\TBD{N} injected); the contract executed \TBD{0 or count} in A3 and \TBD{count} in A2 (H2 \TBD{supported/not supported}). As declared before the study, the input-triggered resource-abuse probe passed behavioral evidence in \TBD{k}/\TBD{m} A3 streams and was contained by the sandbox.

\textbf{Answer to RQ4.} Not established. The contract's enforcement is effective and inexpensive, but this evaluation cannot show that a governed self-evolving harness is useful: \TBD{floor: models rarely self-repair even unguarded / inferior: the contract costs more than the prespecified margin}. We therefore claim an enforcement contribution, a measured cost of conservative admission, and reusable, frozen infrastructure for measuring usefulness as models improve, not improved task performance.

% ---- abstract replacement for the usefulness sentence
% A prespecified 540-stream study of twelve planted-defect tool families shows that the contract executes none of the \TBD{n} violations that original Harness$^2$ executes, but it cannot establish usefulness: \TBD{floor: success is \TBD{A1} even without the contract / inferior: success falls \TBD{A1-A3} points below the unguarded system}.

% ---- framing note (title/venue decision follows, see README)
% The title stays enforcement-focused; claims about self-improvement are removed from the introduction's contribution list and stated as future work.
